\section{CES Representation at Macro Level}
\label{sec:aggregation}

Sections~\ref{sec:shortrun}--\ref{sec:longrun} study a single firm's choices, whereas the automation literature models production directly at the level of the aggregate economy \citep{acemoglu2018, acemoglu2022tasks, acemoglu2025simple}.
This section shows that, under additional assumptions, the two views connect.
A firm's many-input, task-level Leontief technology reduces to a two-input Leontief in skill-adjusted AI management labor and skill-adjusted manual labor (Appendix~\ref{sec:agg_within}), and aggregating across firms that differ only in how effectively they deploy the same AI technology yields an economy-wide CES production function.
This tractably links micro-level organization decisions to the macroeconomic implications of improving AI quality.

Before the aggregation analysis, let us first describe the setup of production.
A firm uses two production inputs: (1) skill-adjusted AI management labor (to complete AI chains), and (2) skill-adjusted manual labor (to perform manual tasks).
These two types of labor demand different compensations, in that executing an AI chain with skill and time cost of one demands $w_{\AIletter}$, while completing a manual task with similar skill and time requirements commands a compensation of $w_{\manualLetter}$.
We call $w_{\AIletter}$ and $w_{\manualLetter}$ the \emph{base wage rates} for AI management and manual labor, respectively.

Firms use these two inputs to complete tasks.
The skill and time requirements of tasks are set by the AI strategy $\T$, and the skill requirement of a job by $\T$ and the job design $\J$ together: a worker's compensation per unit of time depends on all tasks in the job, so the more tasks a job includes, the higher the compensation required for every task in it.
The total compensation of a job is the sum of the skill costs of its tasks, each weighted by the base wage rate of the labor that performs it.
Specifically, the compensation for job $J$ will be:
\begin{equation}
w_{\manualLetter} \left(\sum_{T_b^{\manualLetter} \in J} \manualSkill{b}\right) + w_{\AIletter} \left(\sum_{T_b^{\AIletter} \in J} \AIskill{b}\right),\label{eq:new_wage}
\end{equation}
which is composed of the sum of contributions of skill costs from its manual tasks denoted by set $T_b^{\manualLetter}$ (each weighted by $w_{\manualLetter}$) and its AI chains denoted by set $T_b^{\AIletter}$ (each weighted by $w_{\AIletter}$).\footnote{
This formulation is essentially equivalent to the original definition of wage in \eqref{eq:wage}.
If we multiply each skill cost in Equation~\eqref{eq:new_wage} by its respective base wage rate and redefine the task's skill cost as the resulting product, we arrive at the original formulation given by \eqref{eq:wage}.
The formulation in \eqref{eq:new_wage} allows tasks to contribute differently to the job's wage depending on their execution mode and the type of labor performing them.}

Next, we define skill-adjusted labor, which is the smallest unit of work in our analysis.
Let $J(b)$ denote the job to which task $b$ belongs, and $E(b)$ be the mode of execution of task $b$ (i.e., $E(b) = \manualLetter$ if $b$ is a manual task, or $E(b) = \AIletter$ if $b$ is an AI chain).
Write $d_b$ for the total difficulty of the steps task $b$ hands to AI, so that $d_b = \sum_{i=\ell}^{r} d_i$ for a chain spanning $(s_\ell, \dotsc, s_r)$ and $d_b = 0$ for a manual step.
A task's time requirement for a single attempt inherits its mode from $E(b)$ and is written $\timeCostLetter^{E(b)}_{b}$, which is $\manualTime{b}$ when task $b$ is a manual step, and $\AItime{b}$, the cost of verifying the chain's augmented endpoint, when it is an AI chain.
Recall from Definition~\ref{def:ai_strategy} that the \emph{expected} time cost of a task is $\timecost{b} = \timeCostLetter^{E(b)}_{b}\,\alpha^{-d_b}$, so that below we write the failure inflation $\alpha^{-d_b}$ explicitly wherever an expected quantity is formed rather than folding it into $\timeCostLetter^{E(b)}_{b}$.
Define the skill-adjusted time requirement of task $b$ as
\[
\skillAdjustedTimeLetter_b = \frac{w_{\manualLetter} \left(\sum_{T_{\ell}^{\manualLetter} \in J(b)} \manualSkill{\ell}\right) + w_{\AIletter} \left(\sum_{T_{\ell}^{\AIletter} \in J(b)} \AIskill{\ell}\right)}{w_{E(b)}} \, \timeCostLetter^{E(b)}_{b}.
\]
Note that the fraction appearing behind $\timeCostLetter^{E(b)}_{b}$ represents an effective skill adjustment factor, as the numerator captures the total compensation of the job task $b$ belongs to, while the denominator normalizes by the base wage rate for the specific mode of execution of task $b$.
The resulting fraction thus has units of skill intensity.
This characterization becomes useful later as it allows expressing the (expected) wage bill paid for task $b$ simply as $w_{\AIletter}\,\skillAdjustedTimeLetter_{b}\,\alpha^{-d_b}$ if it is an AI chain or as $w_{\manualLetter}\,\skillAdjustedTimeLetter_{b}$ if it is executed manually.


\subsection{Within-Firm Aggregation: Leontief to Leontief}
\label{sec:agg_within}

We now consider the firm's production.
To produce output, each firm must complete requirements of all \( m \) production steps \( s_1, \ldots, s_m \).
This naturally results in a Leontief production structure with multiple inputs.\footnote{
Different AI deployment strategies and job designs lead to different input proportions but share the same fundamental Leontief form.
}
We show that the firm's production function can be represented by a Leontief function with two firm-level aggregate inputs: AI management labor and manual labor.

Suppose the firm solves the long-run cost minimization problem in \eqref{eq:totalcost_with_handoff} for a given AI quality level $\alpha$, and obtains the optimal AI strategy \( \T \) and job design \( \J \).
Let \( |\T| = n \), implying the original \( m \) production steps are organized into \( n \) distinct tasks, and \( |\J| = p \), indicating these \( n \) tasks are grouped into \( p \) jobs.
With $\T$ and $\J$ fixed, each step's execution mode, the job boundaries, and each job's total skill requirement are given, so we can work directly with tasks rather than steps.

We assume hand-offs are performed manually by the worker who completes the final task of each job.
This allows us to simplify the analysis in two ways:
(1) Since job design is fixed, we know exactly which steps occur at the job boundaries and thus incur hand-off time costs. Therefore, hand-off of job $J_j$ can be treated as a standalone, human-executed task with time requirement \( \handofftime{}(J_j) \) and skill requirement equal to the total skill requirement of job $J_j$.\footnote{
Denote the skill-adjusted time cost of hand-off task of job $J_j$ with $\skillAdjustedTimeHandoff{}(J_j)$.
}
(2) The order of tasks, including the newly defined hand-off tasks, can be rearranged within the production sequence.
Specifically, we relabel tasks such that tasks \( 1 \) to \( k \) are AI chains and thus require AI management labor, while tasks \( k+1 \) to \( n \), along with the hand-off tasks \( n+1 \) to \( n+p-1 \), are executed manually and require manual labor.\footnote{
The final job \( p \) requires no hand-off by definition, so the last hand-off occurs at the end of job \( p-1 \).
}

The firm's task-level production can be expressed in the following Leontief form:
\begin{equation}
y = \min\left\{\frac{\labor{1}}{\skillAdjustedTimeAI{1}\,\alpha^{-d_{1}}},\,\cdots,\,\frac{\labor{k}}{\skillAdjustedTimeAI{k}\,\alpha^{-d_{k}}},\,\frac{\labor{k+1}}{\skillAdjustedTimeManual{k+1}},\,\cdots,\,\frac{\labor{n}}{\skillAdjustedTimeManual{n}},\,\frac{\labor{n+1}}{\skillAdjustedTimeHandoff{}(J_1)},\,\cdots,\,\frac{\labor{n+p-1}}{\skillAdjustedTimeHandoff{}(J_{p-1})}\right\}.\label{eq:task_level_prod}
\end{equation}
The \( y \) on the left-hand side represents the firm's output while $\labor{b}$ is the amount of labor assigned to task $b$. All other variables remain as previously defined.

Since the AI strategy and job design are fixed, the required amount of skill-adjusted time to spend on each task in the denominators of Equation~\eqref{eq:task_level_prod} are fixed, and the firm takes them as given.
In equilibrium, allocation of labor to tasks satisfies
\[
y=\frac{\labor{1}}{\skillAdjustedTimeAI{1}\,\alpha^{-d_{1}}}=\cdots=\frac{\labor{k}}{\skillAdjustedTimeAI{k}\,\alpha^{-d_{k}}}=\frac{\labor{k+1}}{\skillAdjustedTimeManual{k+1}}=\cdots=\frac{\labor{n}}{\skillAdjustedTimeManual{n}}=\frac{\labor{n+1}}{\skillAdjustedTimeHandoff{}(J_1)}=\cdots=\frac{\labor{n+p-1}}{\skillAdjustedTimeHandoff{}(J_{p-1})}.
\]
Notice that the ratio of labor allocated to any two tasks $a$ and $b$ is independent of output level and wage rates (regardless of whether tasks $a$ and $b$ are executed manually or by the help of AI), and solely depends on their relative skill-adjusted time requirements, which are fixed given $\T$ and $\J$.
The fixed ratio of labor inputs implies that the rate of substitution between any pair of tasks is independent of the labor allocated to other tasks:
\begin{equation}
\frac{\partial}{\partial \labor{z}}\left(\frac{\frac{\Delta y}{\Delta \labor{a}}}{\frac{\Delta y}{\Delta \labor{b}}}\right)=0,\qquad\forall z\neq a,b.\label{eq:task_subst_rate_indep}
\end{equation}

A necessary and sufficient condition to aggregate the firm's production function from a Leontief with one input per task into a Leontief with (firm-level) aggregate AI management labor and manual labor is that the rate of substitution between any two tasks within the same aggregate input type is independent of all tasks in the other aggregate input type \citep{leontief1947introduction, fisher1965embodied, felipe2003aggregation}.

The condition expressed in \eqref{eq:task_subst_rate_indep} satisfies the aggregation requirement above.
Thus, the firm's production function can be represented as:
\begin{equation}
y=\min\,\left\{\frac{\bar{\alpha}\,\labor{\AIletter}}{\skillAdjustedTimeLetter_{\AIletter}},\,\frac{\labor{\manualLetter}}{\skillAdjustedTimeLetter_{\manualLetter}}\right\}.\label{eq:micro_agg_prod}
\end{equation}
Here, $\labor{\AIletter}$ and $\labor{\manualLetter}$ represent respectively the firm-level aggregate AI management labor and manual labor.
Moreover, the aggregate skill-adjusted AI and manual time requirements, $\skillAdjustedTimeLetter_{\AIletter}$ and $\skillAdjustedTimeLetter_{\manualLetter}$, are defined as the sum of skill-adjusted time requirements of tasks in their respective groups:
\begin{equation}
\skillAdjustedTimeLetter_{\AIletter} = \sum_{b=1}^{k}\skillAdjustedTimeAI{b}, \qquad
\skillAdjustedTimeLetter_{\manualLetter} = \sum_{j=1}^{p-1}\skillAdjustedTimeHandoff{}(J_j) + \sum_{b=k+1}^{n}\skillAdjustedTimeManual{b}. \label{eq:def_t_AI}
\end{equation}
Finally, $\bar{\alpha}$ is the firm's effective AI quality level defined as
\begin{equation}
\bar{\alpha}=\frac{\sum_{b=1}^{k}\skillAdjustedTimeAI{b}}{\sum_{b=1}^{k}\skillAdjustedTimeAI{b}\alpha^{-d_{b}}},\label{eq:alpha_bar}
\end{equation}
so that the following relationship holds:
\[
\frac{\skillAdjustedTimeLetter_{\AIletter}}{\bar{\alpha}}=\sum_{b=1}^{k}\skillAdjustedTimeAI{b}\alpha^{-d_{b}}.
\]

In short, Equation~\eqref{eq:micro_agg_prod} allows us to analyze labor allocation using the simpler two-input Leontief production function for a given AI strategy \( \T \) and job design \( \J \).
This firm-level aggregate Leontief function implies the equilibrium condition:
\begin{equation}
y=\frac{\bar{\alpha}\,\labor{\AIletter}}{\skillAdjustedTimeLetter_{\AIletter}}
=\,\frac{\labor{\manualLetter}}{\skillAdjustedTimeLetter_{\manualLetter}}.
\label{eq:micro_fixed_proportions}
\end{equation}



\subsection{Cross-Firm Aggregation: Leontief to CES}

Next, we ask whether the Leontief production functions of many firms can be represented by a single economy-wide production function whose inputs aggregate the corresponding firm-level inputs.
The aggregation literature shows that suitable heterogeneity across firms can smooth micro-level production functions into a well-behaved aggregate one (see \cite{sato1975production}, Chapters 2 and 4).\footnote{
For instance, \cite{houthakker1955pareto} studies aggregation from micro Leontief functions to a macro Cobb-Douglas function; \cite{levhari1968note} investigates aggregation from micro Leontief functions to a macro CES function; and \cite{sato1969micro} explores aggregation from micro CES to macro CES functions. 
See \cite{baqaee2019foundations} for a broader formulation of this problem.}
Most of its existence results do not yield a closed form; we show that under certain conditions on firm heterogeneity, firm-level Leontief production functions aggregate to an economy-wide CES production function.

Consider a unit mass of firms in the economy, each producing output with AI management labor, human labor, and capital.
We assume capital is a fixed input whose productivity is common across firms, so that firms make their labor decisions conditional on a given capital stock and capital plays no part in the allocation in \eqref{eq:micro_agg_prod}.\footnote{Since we introduce firm-level heterogeneity through variations in effective AI quality alone, we abstract away from additional sources of heterogeneity across firms.}
We normalize the economy's aggregate capital stock to $1$.
Firms do not know their individual effective AI quality level before production occurs; instead, they share a common belief about the distribution from which these quality levels are drawn.

Production unfolds in two stages.
In the first stage, firms solve the long-run cost minimization problem in \eqref{eq:totalcost_with_handoff}, committing to an AI deployment strategy $\T$ and a job design $\J$ based on their \emph{expected} AI quality levels.
Since firms hold the same expectations about the distribution of effective AI quality, they choose identical AI strategies and job designs.
In the second stage, firms learn their realized effective AI quality levels, hire the required labor and capital, and begin production.\footnote{As production depends on the realized effective AI quality level, $\bar{\alpha}$ also determines firms' sizes.}
The resulting Leontief production function, determined by the previously selected AI strategy, job design, and the firm's realized effective AI quality, takes a form similar to Equation~\eqref{eq:micro_agg_prod}.

In what follows, we demonstrate that the heterogeneity in effective AI quality level can be structured so that micro-level Leontief production functions in \eqref{eq:micro_agg_prod} aggregate into a macro-level CES production function, in which the CES inputs are aggregates of the micro-level inputs \citep{levhari1968note}.
Suppose specifically that the macro-level production function takes the form:
\begin{equation}
Y = \left(\theta_{\AIletter}\,\aggregateAIlabor^{\rho}+\theta_{\manualLetter}\,\aggregateManualLabor^{\rho}+(1-\theta_{\AIletter}-\theta_{\manualLetter})\,K^\rho\right)^{\frac{1}{\rho}},\label{eq:macro_agg_prod}
\end{equation}
where $Y$ represents the (economy-wide) aggregate output, $\aggregateAIlabor$ denotes aggregate AI management labor, $\aggregateManualLabor$ denotes aggregate manual labor, $K$ is aggregate capital, and parameters $\theta_{\AIletter}$ and $\theta_{\manualLetter}$ represent the weights of corresponding inputs.
Aggregate capital is normalized to $1$ as described above, so the third term is constant at $1-\theta_{\AIletter}-\theta_{\manualLetter}$; its exponent is part of the CES form we posit rather than something the aggregation derives.

To link the macro production function above to the micro production functions, we must relate aggregated firm-level inputs \( \labor{\AIletter}\) and \( \labor{\manualLetter} \) from Equation~\eqref{eq:micro_agg_prod} to their corresponding economy-wide aggregates \( \aggregateAIlabor\) and \( \aggregateManualLabor \) in Equation~\eqref{eq:macro_agg_prod}.
This requires either determining the macro production function parameters from the distribution of heterogeneity, or specifying the form of heterogeneity given a set of macro production function parameters.
We adopt the latter approach and derive the output density function in terms of effective AI quality, denoted by \( \phi(\bar{\alpha}) \), as a function of \( \theta_{\AIletter}, \theta_{\manualLetter}, \rho \).
Throughout, we assume $\rho<0$, so the elasticity of substitution is $\sigma<1$, as is usual in this literature.

Normalize the output price to $p=1$, so that $w_{\AIletter}$ and $w_{\manualLetter}$ can be interpreted as real wage rates for a unit of skill-adjusted AI management and manual labor, respectively.
A firm produces only if it earns nonnegative profits:
\[
w_{\AIletter}\,\labor{\AIletter}+w_{\manualLetter}\,\labor{\manualLetter}\leq y.
\]
Substituting for $\labor{\manualLetter}$ and $y$ from Condition~\eqref{eq:micro_fixed_proportions} into the profitability condition yields:
\[
w_{\AIletter}\,\labor{\AIletter}+w_{\manualLetter}\,\frac{\skillAdjustedTimeLetter_{\manualLetter}\,\bar{\alpha}\,\labor{\AIletter}}{\skillAdjustedTimeLetter_{\AIletter}} \leq \frac{\bar{\alpha}\,\labor{\AIletter}}{\skillAdjustedTimeLetter_{\AIletter}}.
\]
Rearranging terms and canceling $\labor{\AIletter}$ gives the lower bound of firms' effective AI quality level:\footnote{We assume the parameters $w_{\AIletter},\,w_{\manualLetter},\,\skillAdjustedTimeLetter_{\AIletter},\,\skillAdjustedTimeLetter_{\manualLetter}$ are such that $0 < ({w_{\AIletter}\,\skillAdjustedTimeLetter_{\AIletter}})/(1-w_{\manualLetter}\,\skillAdjustedTimeLetter_{\manualLetter}) < 1$.}
\[
\bar{\alpha} \geq \frac{w_{\AIletter}\,\skillAdjustedTimeLetter_{\AIletter}}{1 - w_{\manualLetter}\,\skillAdjustedTimeLetter_{\manualLetter}}.
\]
This lower bound implies that only firms with realized effective AI quality level above this threshold will produce in equilibrium and others exit the market.
Moreover, observe from \eqref{eq:alpha_bar} that for $d_b \geq 0$ the upper bound of $\bar{\alpha}$ is 1 and is achieved when $\alpha \rightarrow 1^{-}$.
Therefore:
\begin{equation}
\underline{\alpha} \equiv \frac{w_{\AIletter}\,\skillAdjustedTimeLetter_{\AIletter}}{1 - w_{\manualLetter}\,\skillAdjustedTimeLetter_{\manualLetter}} \leq \bar{\alpha} \leq 1.
\label{eq:alpha_threshold}
\end{equation}

Let $\phi(\bar{\alpha})$ denote output across firms by effective AI quality, normalizing the density of firms over $\bar{\alpha}$ to one, so that a firm of type $\bar{\alpha}$ produces $y = \phi(\bar{\alpha})$ and aggregate output is the integral of $\phi$ over the producing firms.
We refer to $\phi$ as the distribution of effective AI quality in this sense.
From Condition~\eqref{eq:micro_fixed_proportions}, it directly follows that the AI management labor and manual labor used by this firm are: $\labor{\AIletter} = (\skillAdjustedTimeLetter_{\AIletter}/\bar{\alpha})\,\phi(\bar{\alpha}),$ and $\labor{\manualLetter} = \skillAdjustedTimeLetter_{\manualLetter}\,\phi(\bar{\alpha}).$
Consequently, aggregate output $Y$, aggregate AI management labor $\aggregateAIlabor$, and aggregate manual labor $\aggregateManualLabor$ can be expressed as:
\begin{equation}
Y = \int_{\underline{\alpha}}^{1}\phi(\bar{\alpha})\,d\bar{\alpha}, \qquad
\aggregateAIlabor = \int_{\underline{\alpha}}^{1}\skillAdjustedTimeLetter_{\AIletter}\frac{\phi(\bar{\alpha})}{\bar{\alpha}}\,d\bar{\alpha}, \qquad
\aggregateManualLabor = \int_{\underline{\alpha}}^{1}\skillAdjustedTimeLetter_{\manualLetter}\,\phi(\bar{\alpha})\,d\bar{\alpha}.
\label{eq:micro_Y}
\end{equation}

% sloppypar: four \eqref tags in one sentence, and under the OA./SA. numbering they are
% "(OA.C.11)" rather than "(C.11)", which is enough extra width that TeX could not
% break the line within the house \tolerance and ran 18pt into the right margin.
\begin{sloppypar}
Substituting the aggregated variables from the micro-level Equation~\eqref{eq:micro_Y} into the aggregate production function in \eqref{eq:macro_agg_prod}, we obtain:
\end{sloppypar}
% \small for the same reason as the display below it: at full size this one is 505pt
% wide against a 470pt measure and its closing exponent ran into the right margin.
% It needs no line break, since \small brings it to 464pt on its own.
\begingroup\smalldisplay
\begin{equation}
\int_{\underline{\alpha}}^{1}\,\phi(\bar{\alpha})\,d\bar{\alpha}
=
\left(
\theta_{\AIletter}\left(\skillAdjustedTimeLetter_{\AIletter}\int_{\underline{\alpha}}^{1}\,\frac{\phi(\bar{\alpha})}{\bar{\alpha}}\,d\bar{\alpha}\right)^{\rho}
+
\theta_{\manualLetter} \left(\skillAdjustedTimeLetter_{\manualLetter}\int_{\underline{\alpha}}^{1}\,\phi(\bar{\alpha})\,d\bar{\alpha}\right)^{\rho}
+
\left(1-\theta_{\AIletter}-\theta_{\manualLetter}\right)
\right)^{\frac{1}{\rho}},
\label{eq:macro_agg_prod_with_micro_aggregates}
\end{equation}
\endgroup
provided that the CES weights and the skill-adjusted time requirements satisfy
\begin{equation}
\theta_{\AIletter}\,\skillAdjustedTimeLetter_{\AIletter}^{\rho}+\theta_{\manualLetter}\,\skillAdjustedTimeLetter_{\manualLetter}^{\rho}=1,
\label{eq:ces_share_restriction}
\end{equation}
which is what the boundary condition at $\bar{\alpha} \to 1$ imposes.

Solving for $\phi(\bar{\alpha})$ in terms of $\theta_{\AIletter}$, $\theta_{\manualLetter}$, and $\rho$ then gives the following distribution for effective AI quality:
\begin{equation}
\phi(\bar{\alpha})=\frac{1}{1-\rho}\left(\frac{1-\theta_{\AIletter}-\theta_{\manualLetter}}{\theta_{\AIletter}\,\skillAdjustedTimeLetter_{\AIletter}^{\rho}}\right)^{\frac{1}{\rho}}(\bar{\alpha})^{\frac{1}{\rho-1}}\left[1-(\bar{\alpha})^{\frac{\rho}{\rho-1}}\right]^{-\frac{1+\rho}{\rho}}.
\label{eq:phi}
\end{equation}
Subsection~\ref{sec:aggregationDist} below derives this distribution.

This completes the production function aggregation procedure.


Two features of this construction are worth making explicit.
First, because $\skillAdjustedTimeLetter_{\manualLetter}$ is common to all firms, Equation~\eqref{eq:micro_Y} gives $\aggregateManualLabor = \skillAdjustedTimeLetter_{\manualLetter}\,Y$ at every wage vector, while capital is normalized to $K=1$.
As the participation threshold varies, the economy therefore traces a one-dimensional locus along which $\aggregateManualLabor/Y$ and $K$ are both fixed, rather than filling out the space of input combinations.
Equation~\eqref{eq:macro_agg_prod} should accordingly be read as a representation valid along that locus rather than as an identification of the economy's technology over arbitrary $(\aggregateAIlabor, \aggregateManualLabor, K)$.
Second, recovering a genuine three-input CES requires firms to differ in at least two of their input requirements, as in \citet{houthakker1955pareto} and \citet{levhari1968note}.
Our two-stage timing rules this out by construction, since every firm commits to the same $\T$ and $\J$ before learning $\bar{\alpha}$.
Allowing firms that anticipate deploying AI more effectively to organize work differently would make $\skillAdjustedTimeLetter_{\manualLetter}$ vary with $\bar{\alpha}$ and restore the missing dimension, at the cost of a joint distribution over organizational choices and realized quality that would be intractable without imposing further assumptions on the type of heterogeneity and firm behavior.





\subsection{Derivation of the Effective AI Quality Heterogeneity Distribution}
\label{sec:aggregationDist}

We derive \eqref{eq:phi} by extending the method of \cite{levhari1968note}.

Define $u=({w_{\AIletter}\,\skillAdjustedTimeLetter_{\AIletter}})/(1-w_{\manualLetter}\,\skillAdjustedTimeLetter_{\manualLetter})$ and rewrite Equation~\eqref{eq:macro_agg_prod_with_micro_aggregates} as:
\begin{equation}
\left(\int_{u}^{1}\,\phi(\bar{\alpha})\,d\bar{\alpha}\right)^{\rho}
=
\theta_{\AIletter}\left(\skillAdjustedTimeLetter_{\AIletter}\int_{u}^{1}\,\frac{\phi(\bar{\alpha})}{\bar{\alpha}}\,d\bar{\alpha}\right)^{\rho}
+
\theta_{\manualLetter} \left(\skillAdjustedTimeLetter_{\manualLetter}\int_{u}^{1}\,\phi(\bar{\alpha})\,d\bar{\alpha}\right)^{\rho}
+
\left(1-\theta_{\AIletter}-\theta_{\manualLetter}\right).
\label{eq:macro_agg_prod_with_hetDist}
\end{equation}
Define:
\begin{equation}
\Gamma(u)=\int_{u}^{1}\,\phi(\bar{\alpha})\,d\bar{\alpha},\label{eq:gamma}
\end{equation}
and
\begin{equation}
\Psi(u)=\int_{u}^{1}\,\frac{\phi(\bar{\alpha})}{\bar{\alpha}}\,d\bar{\alpha}.\label{eq:psi}
\end{equation}
Note that $\Gamma^{'}(u)=-\phi(u)$ and $\Psi^{'}(u)=-\phi(u)/u$.
Rewrite Equation~\eqref{eq:macro_agg_prod_with_hetDist} as:
\[
\Gamma^{\rho}(u)=\theta_{\AIletter}\,\skillAdjustedTimeLetter_{\AIletter}^{\rho}\,\Psi^{\rho}(u)+\theta_{\manualLetter}\,\skillAdjustedTimeLetter_{\manualLetter}^{\rho}\,\Gamma^{\rho}(u)+(1-\theta_{\AIletter}-\theta_{\manualLetter}).
\]
Rearranging terms, we obtain:
\begin{equation}
\left(1-\theta_{\manualLetter}\,\skillAdjustedTimeLetter_{\manualLetter}^{\rho}\right)\Gamma^{\rho}(u)=\theta_{\AIletter}\,\skillAdjustedTimeLetter_{\AIletter}^{\rho}\,\Psi^{\rho}(u)+(1-\theta_{\AIletter}-\theta_{\manualLetter}).\label{eq:macro_agg_pro_with_gamma}
\end{equation}

Differentiating Equation~\eqref{eq:macro_agg_pro_with_gamma} with respect to $u$ yields:
\[
\left(1-\theta_{\manualLetter}\,\skillAdjustedTimeLetter_{\manualLetter}^{\rho}\right)\Gamma^{\rho-1}(u)=\theta_{\AIletter}\,\skillAdjustedTimeLetter_{\AIletter}^{\rho}\,\Psi^{\rho-1}(u)\,u^{-1}.
\]
Solving for $\Psi^{\rho-1}(u)$ as a function of $u$ and $\Gamma^{\rho-1}(u)$ we get:
\[
\Psi^{\rho-1}(u)=\frac{\left(1-\theta_{\manualLetter}\,\skillAdjustedTimeLetter_{\manualLetter}^{\rho}\right)}{\theta_{\AIletter}\,\skillAdjustedTimeLetter_{\AIletter}^{\rho}}\Gamma^{\rho-1}(u)\,u.
\]
Finally, express $\Psi(u)$ in terms of $\Gamma(u)$:
\begin{equation}
\Psi(u)=\left(1-\theta_{\manualLetter}\,\skillAdjustedTimeLetter_{\manualLetter}^{\rho}\right)^{\frac{1}{\rho-1}}\left(\theta_{\AIletter}\,\skillAdjustedTimeLetter_{\AIletter}^{\rho}\right)^{\frac{1}{1-\rho}}u^{\frac{1}{\rho-1}}\,\Gamma(u).\label{eq:psi_solved}
\end{equation}


Substituting \( \Psi(u) \) from Equation~\eqref{eq:psi_solved} into Equation~\eqref{eq:macro_agg_pro_with_gamma}, we have:
\[
\left(1-\theta_{\manualLetter}\,\skillAdjustedTimeLetter_{\manualLetter}^{\rho}\right)\Gamma^{\rho}(u)=\left(1-\theta_{\manualLetter}\,\skillAdjustedTimeLetter_{\manualLetter}^{\rho}\right)^{\frac{\rho}{\rho-1}}\left(\theta_{\AIletter}\,\skillAdjustedTimeLetter_{\AIletter}^{\rho}\right)^{\frac{1}{1-\rho}}u^{\frac{\rho}{\rho-1}}\,\Gamma^{\rho}(u)+\left(1-\theta_{\AIletter}-\theta_{\manualLetter}\right).
\]
Rearranging terms to solve for \( \Gamma(u) \), we get:
\[
\Gamma(u)=\left(1-\theta_{\AIletter}-\theta_{\manualLetter}\right)^{\frac{1}{\rho}}\left[1-\theta_{\manualLetter}\,\skillAdjustedTimeLetter_{\manualLetter}^{\rho}-\left(1-\theta_{\manualLetter}\,\skillAdjustedTimeLetter_{\manualLetter}^{\rho}\right)^{\frac{\rho}{\rho-1}}\left(\theta_{\AIletter}\,\skillAdjustedTimeLetter_{\AIletter}^{\rho}\right)^{\frac{1}{1-\rho}}u^{\frac{\rho}{\rho-1}}\right]^{-\frac{1}{\rho}}.
\]

The relation $\Gamma'=-\phi$ pins down $\Gamma$ only up to an additive constant, so the expression above must still be checked against the boundary condition $\Gamma(1)=0$ that \eqref{eq:gamma} imposes by construction.
Evaluating the bracket at $u=1$ leaves
\[
1-\theta_{\manualLetter}\,\skillAdjustedTimeLetter_{\manualLetter}^{\rho}-\left(1-\theta_{\manualLetter}\,\skillAdjustedTimeLetter_{\manualLetter}^{\rho}\right)^{\frac{\rho}{\rho-1}}\left(\theta_{\AIletter}\,\skillAdjustedTimeLetter_{\AIletter}^{\rho}\right)^{\frac{1}{1-\rho}},
\]
which vanishes if and only if $1-\theta_{\manualLetter}\,\skillAdjustedTimeLetter_{\manualLetter}^{\rho}=\theta_{\AIletter}\,\skillAdjustedTimeLetter_{\AIletter}^{\rho}$, that is, if and only if Restriction~\eqref{eq:ces_share_restriction} holds.
Under that restriction the bracket collapses to $\theta_{\AIletter}\,\skillAdjustedTimeLetter_{\AIletter}^{\rho}\bigl[1-u^{\rho/(\rho-1)}\bigr]$, so that
\begin{equation}
\Gamma(u)=\left(\frac{1-\theta_{\AIletter}-\theta_{\manualLetter}}{\theta_{\AIletter}\,\skillAdjustedTimeLetter_{\AIletter}^{\rho}}\right)^{\frac{1}{\rho}}\left[1-u^{\frac{\rho}{\rho-1}}\right]^{-\frac{1}{\rho}},
\label{eq:gamma_solved}
\end{equation}
which satisfies $\Gamma(1)=0$ because $-1/\rho>0$ for $\rho<0$.
Restriction~\eqref{eq:ces_share_restriction} is therefore not an auxiliary assumption but the condition under which a density consistent with \eqref{eq:gamma} exists at all.
It ties the CES weights to $\skillAdjustedTimeLetter_{\AIletter}$ and $\skillAdjustedTimeLetter_{\manualLetter}$, and hence to the firm's AI strategy and job design.
It also delivers $1-\theta_{\manualLetter}\,\skillAdjustedTimeLetter_{\manualLetter}^{\rho}=\theta_{\AIletter}\,\skillAdjustedTimeLetter_{\AIletter}^{\rho}>0$, so the real powers appearing throughout are well defined.

Recall that \( \Gamma'(u)=-\phi(u) \).
Differentiating \eqref{eq:gamma_solved}, we obtain:
\[
\Gamma'(u)=\frac{1}{\rho-1}\left(\frac{1-\theta_{\AIletter}-\theta_{\manualLetter}}{\theta_{\AIletter}\,\skillAdjustedTimeLetter_{\AIletter}^{\rho}}\right)^{\frac{1}{\rho}}u^{\frac{1}{\rho-1}}\left[1-u^{\frac{\rho}{\rho-1}}\right]^{-\frac{1+\rho}{\rho}}.
\]
Thus, we have:
\[
\phi(\bar{\alpha})=\frac{1}{1-\rho}\left(\frac{1-\theta_{\AIletter}-\theta_{\manualLetter}}{\theta_{\AIletter}\,\skillAdjustedTimeLetter_{\AIletter}^{\rho}}\right)^{\frac{1}{\rho}}(\bar{\alpha})^{\frac{1}{\rho-1}}\left[1-(\bar{\alpha})^{\frac{\rho}{\rho-1}}\right]^{-\frac{1+\rho}{\rho}},
\]
which is the formula provided in \eqref{eq:phi} above.
